\documentclass[10pt,a4paper]{amsart}
\usepackage[margin=19mm]{geometry}
\usepackage{amsmath,amssymb,mathtools}
\usepackage[scr=rsfs,cal=euler]{mathalfa}
\usepackage{xfrac}
\usepackage[unicode,hidelinks,bookmarks=false]{hyperref}
\newcommand{\ie}{i.e.\ }
\newcommand{\cf}{cf.\ }
\newcommand{\resp}{respectively\ }
\newcommand{\Subsection}{\subsection}
\providecommand{\nobreakdash}{\nobreak\hbox{-}}
\setlength{\emergencystretch}{2em}
\multlinegap0pt
% ---------------------------------------------------------------------------
% Editorial AMS wrapper prelude. The theorem environments and the mathematical macros below are
% transcribed from the paper's own preamble (cached-originals/source0.tex lines 60--195); the AMS amsart
% class, the named format packages of the pinned TeX Live runtime and this editorial formatting are not
% part of the author's text. No mathematics is changed.
% ---------------------------------------------------------------------------




\newcommand{\area}{\mathsf{area}}
\newcommand{\dinv}{\mathsf{dinv}}
\newcommand{\Dinv}{\mathsf{Dinv}}
\newcommand{\inv}{\mathsf{inv}}
\newcommand{\Inv}{\mathsf{Inv}}
\newcommand{\Des}{\mathsf{Des}}
\newcommand{\des}{\mathsf{des}}
\newcommand{\lec}{\mathsf{lec}}
\newcommand{\single}{\mathsf{single}}
\newcommand{\pix}{\mathsf{pix}}
\newcommand{\sw}{\mathsf{sw}}
\newcommand{\srw}{\mathsf{srw}}
\newcommand{\ins}{\mathsf{ins}}
\newcommand{\rins}{\mathsf{rins}}
\newcommand{\h}{\mathsf{h}}
\newcommand{\wv}{\mathsf{wv}}
\newcommand{\wave}{\mathsf{wave}}
\newcommand{\rwave}{\mathsf{rwave}}
\newcommand{\conc}{\mathsf{conc}}
\newcommand{\srdes}{\mathsf{srdes}}

\newcommand{\pmaj}{\mathsf{pmaj}}
\newcommand{\maj}{\mathsf{maj}}
\newcommand{\ides}{\mathsf{ides}}



\newcommand{\dcomp}{\mathsf{dcomp}}
\newcommand{\comp}{\mathsf{comp}}
\newcommand{\st}{\mathsf{st}}
\newcommand{\set}{\mathsf{set}}
\newcommand{\asc}{\mathsf{asc}}
\newcommand{\LLT}{\mathrm{LLT}}



\newcommand{\D}{\mathsf{D}} % Dyck paths (new)
\newcommand{\W}{\mathsf{W}} % Labellings

\newcommand{\LD}{\mathsf{LD}} % labelled Dyck paths


\newcommand{\ISF}{\mathsf{ISF}}
\newcommand{\wt}{\mathsf{wt}}
\newcommand{\op}{\mathsf{op}}



\newcommand{\Hess}{\mathcal Hess}
\newcommand{\cN}{\mathcal N}
\newcommand{\cM}{\mathcal M}
\newcommand{\cL}{\mathcal L}
\newcommand{\cX}{\mathcal X}

\newcommand{\bN}{\mathbb N}
\newcommand{\bZ}{\mathbb Z}
\newcommand{\bP}{\mathbb P}
\newcommand{\bS}{\mathbf {S}}

\newcommand{\e}{\underline{e}}
\newcommand{\PPi}{\mathbf{\Pi}}


\newcommand{\qbinom}[2]{\genfrac{[}{]}{0pt}{}{#1}{#2}}
\newcommand{\dqbinom}[2]{\genfrac{[}{]}{0pt}{0}{#1}{#2}}



\newtheorem{theorem}{Theorem}[section]
\newtheorem{lemma}[theorem]{Lemma}
\newtheorem{proposition}[theorem]{Proposition}
\newtheorem{corollary}[theorem]{Corollary}
\newtheorem{conjecture}[theorem]{Conjecture}

\newtheorem{manualtheoreminner}{Theorem}
\newenvironment{manualtheorem}[1]{%
	\renewcommand\themanualtheoreminner{#1}%
	\manualtheoreminner
}{\endmanualtheoreminner}

\newtheorem{manualconjinner}{Conjecture}
\newenvironment{manualconj}[1]{%
	\renewcommand\themanualconjinner{#1}%
	\manualconjinner
}{\endmanualconjinner}

\theoremstyle{definition}
\newtheorem{definition}[theorem]{Definition}
\newtheorem{definitions}[theorem]{Definitions}
\newtheorem{example}[theorem]{Example}
\newtheorem{xca}[theorem]{Exercise}

\theoremstyle{remark}
\newtheorem{remark}[theorem]{Remark}

\numberwithin{equation}{section}



\begin{document}
\title{Generalizing Eulerian Numbers via Semipermutations: Topological and Combinatorial Aspects}
\author{Giovanni Gaiffi and Giovanni Interdonato}
\date{Source: arXiv:2601.17945v2; distributed under CC BY 4.0}
\maketitle
\noindent\emph{Source, version and licence.} \emph{Generalizing Eulerian Numbers via Semipermutations: Topological and Combinatorial Aspects}, by Giovanni Gaiffi and Giovanni Interdonato. The source of this editable unit is the e-print of \textbf{version v2} of arXiv:2601.17945, https://arxiv.org/abs/2601.17945v2, whose material is distributed under the Creative Commons Attribution 4.0 International licence, http://creativecommons.org/licenses/by/4.0/, as recorded in the arXiv licence metadata for this paper and shown on that abstract page. The whole of the body below is version v2, the newest version of the paper. Full rights notice: licenses/arxiv-2601.17945-CC-BY-4.0.txt.\par\smallskip
\noindent\textit{Editorial scope.} Counted unit: original Theorem 5.11 of the article, the statement carried by the source environment \texttt{theorem} with source label \texttt{thm:bijection}, cached-originals/source0.tex lines 838--844, printed as \textbf{Theorem 5.11} exactly as in the article. It is delivered together with the authors' own complete proof as the single primary proof body, source0.tex lines 846--892 (47 lines): the \texttt{proof} environment that immediately follows the statement, runs the induction on the length of the semipermutation, proves that $\psi$ is well defined and bijective, and identifies the two statistics. No proof marker is added and no mathematics is written, repaired or re-derived; the authors' notation and the printed slips of the source are kept literal. The closure of the counted statement and proof over references and citations was computed: it contains no citation at all, and the definition of the map $\psi$ together with the auxiliary statements the proof invokes are included in the delivered file as uncounted context, so every reference inside the retained passages resolves inside it. The article's own numbering is reproduced by explicit counter settings: the counter \texttt{theorem}, declared by \texttt{\textbackslash newtheorem\{theorem\}\{Theorem\}[section]} with lemma, proposition, corollary, conjecture, definition, example and exercise sharing it, is set before each retained passage, so the counted statement prints as Theorem 5.11 exactly as in the article. No \texttt{\textbackslash cite} command occurs in any retained passage, so this unit delivers no bibliography entry; the article's own bibliography data file \texttt{Biblebib.bib} is neither spliced into the bundled source nor redistributed. The retained passages contain no figure environment and no graphics-inclusion command, so no artwork is redistributed. The source's own class is AMS \texttt{amsart} (\texttt{\textbackslash documentclass[a4paper]\{amsart\}}); the e-print ships no class or style file, so no publisher class is shipped, used or redistributed, and the wrapper is the same AMS \texttt{amsart} class with the article's own macros and theorem declarations transcribed from its preamble. Every declared body of this unit is an exact, unmodified span of source0.tex: no body is a bounded passage comparison, \texttt{omit} was not used anywhere, and consequently no fidelity receipt is asserted in delivery-evidence.json. The complete concatenated original source is bundled as sources/193353/source0.tex. Omitted from this editable unit and therefore absent from it are source0.tex lines 1--277; 281--383; 392--406; 410--755; 767--779; 786--821; 825--837; 845--845; 893--956, that is the article's own preamble, title and author block and abstract, the introduction, the topological background of the configuration spaces, the earlier bijections and their proofs, and the article's bibliography data, all of which remain present in the bundled source but are not delivered here. The AMS \texttt{amsart} wrapper, the named format packages of the pinned TeX Live runtime and this editorial source, licence and scope paragraph are editorial formatting only.\par\smallskip
\setcounter{section}{1}
\setcounter{equation}{0}
\setcounter{manualconjinner}{0}
\setcounter{manualtheoreminner}{0}
\setcounter{theorem}{1}
\begin{remark}\label{ind_lec}
We can obtain this factorization recursively, taking as $\alpha_k$ the suffix of $w$ that starts at the position of the rightmost descent of $w$, and then iterating this reasoning on the remaining prefix obtained by removing $\alpha_k$. The process ends when no further descents are present, i.e. when we have a word whose elements are in increasing order: this will be $\gamma$.
\end{remark}

\setcounter{section}{3}
\setcounter{equation}{0}
\setcounter{manualconjinner}{0}
\setcounter{manualtheoreminner}{0}
\setcounter{theorem}{0}
\begin{remark} \label{lec_inverse}
We can evaluate the $\lec$ of a word $w$ in a recursive way, as suggested by the Remark~\ref{ind_lec}. If $w=\gamma*\alpha_1*\cdots*\alpha_k$ is the Hook factorization, and $k=0$, then $\lec(w)=\lec(\gamma)=0$, while if $k>0$, then we have by definition that
\begin{equation}
    \lec(w)=\lec(\gamma*\alpha_1*\cdots*\alpha_{k-1})+\lec(\alpha_k).
\end{equation}

If a hook $\alpha$ has length $\ell$, then we obtain that $1\leq\lec(\alpha)\leq\ell-1$. Moreover, if we have $\ell$ distinct numbers $a_1<a_2<\cdots<a_{\ell}$, for any $0\leq d\leq\ell-1$, we have exactly one way to rearrange them to create a hook, possibly trivial if $d=0$,  $\alpha_d$ with $\lec(\alpha_d)$ equal to $d$, that is, $\alpha_d=a_{d+1}a_1a_2\cdots a_da_{d+2}\cdots a_{\ell}$.
\end{remark}

\setcounter{section}{3}
\setcounter{equation}{2}
\setcounter{manualconjinner}{0}
\setcounter{manualtheoreminner}{0}
\setcounter{theorem}{3}
\begin{remark}\label{rem:unique_rwave}
Similarly to what was observed in Remark~\ref{lec_inverse}, if a reverse wave $rw$ has length $\ell $, then we have that $1\leq\des(rw)\leq\ell-1$. Moreover if we have $\ell$ distinct numbers $b_1<b_2<\cdots<b_{\ell}$, for any $0\leq d\leq\ell-1$, we have exactly one way to rearrange them to create a reverse wave, possibly trivial if $d=0$, $r_d$ with $\des(rw_d)$ equal to $d$, that is, ${rw_d=b_1b_2\cdots b_{\ell-d-1}b_{\ell}b_{\ell-1}\cdots b_{\ell-d}}$.
\end{remark}

\setcounter{section}{5}
\setcounter{equation}{4}
\setcounter{manualconjinner}{0}
\setcounter{manualtheoreminner}{0}
\setcounter{theorem}{5}
\begin{proposition}\label{prop:inv}
Let $a=a_1\cdots a_p$ and $b=b_1\cdots b_q$ two words in $W$ with distinct entries and such that $b$ is a non trivial wave. Let $0\leq x\leq p$ be the maximum such that $a_1<a_2<\cdots<a_x<b_q$, where $x=0$ if $a_1>b_q$. Then the map $\ins(a,b)=a_1\cdots a_xb_1\cdots b_qa_{x+1}\cdots a_p$ satisfies
\begin{align}\label{eq:ins1}
    \sw(\ins(a,b))=b,\\
    \ins(a,b)\setminus\sw(\ins(a,b))&=a,\label{eq:ins2}
\end{align}
and, for all $w\in W$ that are not increasing,
\begin{equation}\label{eq:ins3}
    \ins(w\setminus\sw(w),\sw(w))=w.
\end{equation}
\end{proposition}

\setcounter{section}{5}
\setcounter{equation}{7}
\setcounter{manualconjinner}{0}
\setcounter{manualtheoreminner}{0}
\setcounter{theorem}{7}
\begin{proposition}\label{prop:des_spezza}
Let $w=w_1\cdots w_{\ell}\in W$ be a word and $\sw(w)=w_rw_{r+1}\cdots w_t\cdots w_{s-1}w_s$ its special wave. Then we have that
\begin{equation}\label{eq:des_spezza}
\des(w)=\des(w\setminus\sw(w))+\des(\sw(w)).
\end{equation}   
\end{proposition}

\setcounter{section}{5}
\setcounter{equation}{9}
\setcounter{manualconjinner}{0}
\setcounter{manualtheoreminner}{0}
\setcounter{theorem}{9}
\begin{remark}\label{rem:unique_wave}
Similarly to what was observed in Remarks~\ref{lec_inverse} and \ref{rem:unique_rwave}, if a wave $\beta$ has length $\ell $, then we have that $1\leq\des(\beta)\leq\ell-1$. Moreover if we have $\ell$ distinct numbers $b_1<b_2<\cdots<b_{\ell}$, for any $0\leq d\leq\ell-1$, we have exactly one way to rearrange them to create a wave, possibly trivial if $d=0$, $\beta_d$ with $\des(\beta_d)$ equal to $d$, that is, ${\beta_d=b_{d+1}b_{d+2}\cdots b_{\ell}b_db_{d-1}\cdots b_1}$.
\end{remark}

\setcounter{section}{5}
\setcounter{equation}{10}
\setcounter{manualconjinner}{0}
\setcounter{manualtheoreminner}{0}
\setcounter{theorem}{10}
\begin{theorem}\label{thm:bijection}
    The map $\psi:W\to W$ just defined is a well defined bijection of $W$, such that \begin{equation}
        \des(w)=\lec(\psi(w)).\label{eq:des_lec}
        \end{equation}

    Moreover the restriction to the symmetric group $\psi_{\mathcal{S}_n}:\mathcal{S}_n\to \mathcal{S}_n$ is a bijection.
\end{theorem}

\setcounter{section}{5}
\setcounter{equation}{11}
\setcounter{manualconjinner}{0}
\setcounter{manualtheoreminner}{0}
\setcounter{theorem}{11}
\begin{proof}
We will prove by induction on the length $\ell$ of $w\in W$ that $\psi(w)$ is a rearrangement of $w$, and that it is a well defined map. If $\ell=0$ then $w=\varepsilon$ is the empty word and by definition $\psi(\varepsilon)=\varepsilon$, that is fine. If $w$ has length $\ell>0$, then by definition $w\setminus\sw(w)$ is a word of length strictly less than $\ell$, and so by the inductive hypothesis $\psi(w\setminus\sw(w))$ is a rearrangement of $w\setminus\sw(w)$. Moreover $\delta(\sw(w))$ is defined as a rearrangement of $\sw(w)$, so this means that $\psi(w)$ is a rearrangement of $w$, and $\psi$ is a well defined map.

To prove that $\psi$ is a bijection we will define its inverse $\varphi$ as the map such that:
\begin{itemize}
    \item If $w=\varepsilon$ is the empty word then $\varphi(\varepsilon)=\varepsilon$.
    \item Otherwise, if the rightmost hook $\h(w)$ of $w$ is not trivial
    \begin{equation}
        \varphi(w)=\ins(\varphi(w\setminus \h(w)),\wv(\h(w))),
    \end{equation}
    where $w\setminus \h(w)$ is the prefix of $w$ obtained removing $\h(w)$ and $$\wv(\h(w))=c_{d+1}c_{d+2}\cdots c_mc_dc_{d-1}\cdots c_1$$ is a rearrangement of $\h(w)$, where $c_1<c_2<\cdots<c_m$ are the elements of $\h(w)$ and $d=\lec(\h(w))=\inv(\h(w))$. If $\h(w)$ is a trivial hook, then $\h(w)=w$ and in this case $\varphi(w)=w$
\end{itemize}
First of all let us note that if $\alpha$ is a hook with $\lec(\alpha)=d$, then, thanks to Remark~\ref{rem:unique_wave}, we have that $\wv(\alpha)$ is the unique wave such that its elements are the same of $\alpha$ and such that $\des(\wv(\alpha))=\lec(\alpha)=d$. Similarly, if $\beta$ is a wave with $\des(\beta)=d$, then, thanks to Remark~\ref{lec_inverse}, we have that $\delta(\beta)$ is the unique hook such that its elements are the same of $\beta$ and such that $\lec(\delta(\beta))=\des(\beta)=d$. Putting these two last observations together, we obtain, if $\alpha$ is a hook and $\beta$ is a wave, that
\begin{align}
    \delta(\wv(\alpha))&=\alpha\label{eq:delta_wv}\\
    \wv(\delta(\beta))&=\beta.\label{eq:wv_delta}
\end{align}


We can notice that from the definition of $\psi$ we have that if $w$ is increasing then $\sw(w)=w$ and $\psi(w)=w$, and so it is clear that $\psi$ and $\varphi$ are inverse when restricted to the subset of increasing words, because they are both the identity in this case.

Now let us consider the case where $w\in W$ is not increasing.

Since $w\setminus\h(w)$ has length strictly less than $w$ and, by definition, $\wv(\h(w))$ is a wave, using inductive reasoning, entirely analogous to what was done previously for $\psi$, we obtain that $\varphi(w)$ is a rearrangement of $w$ and that it is a well defined map. 

Let us prove that they are inverse maps by induction on the length $\ell$ of $w$. If $\ell=0$ it is clear. If $\ell>0$ then
\begin{align}
\psi(\varphi(w))&=\psi(\ins(\varphi(w\setminus\h(w)),\wv(\h(w))))\notag\\
&=\psi(\varphi(w\setminus\h(w)))*\delta(\wv(\h(w)))\\
&=w\setminus\h(w)*\h(w)=w,\notag
\end{align}
where the second equality is true thanks to Proposition~\ref{prop:inv} and the last equality is true thanks to inductive hypothesis, since $w\setminus\h(w)$ has length strictly less than $\ell$, and thanks to equation~\eqref{eq:delta_wv}. Moreover
\begin{align}
    \varphi(\psi(w))&=\ins(\varphi(\psi(w)\setminus\h(\psi(w))),\wv(\h(\psi(w))))\notag\\
    &=\ins(\varphi(\psi(w)\setminus\delta(\sw(w))),\wv(\delta(\sw(w))))\\
    &=\ins(\varphi(\psi(w\setminus\sw(w))),\sw(w))\notag\\
    &=\ins(w\setminus\sw(w)),\sw(w))=w\notag,
\end{align}
where the second equality is true since the rightmost hook of $\psi(w)$ is by definition $\delta(\sw(w))$, the third equality is true thanks to equation~\eqref{eq:wv_delta} and thanks to the definition of $\psi(w)$. The fourth equality holds by inductive hypothesis, since the length of $w\setminus\sw(w)$ is strictly less than $\ell$ and the last is true thanks to Proposition~\ref{prop:inv}. We have proved that $\psi$ and $\varphi$ are inverse, and so that they are bijections of $W$. Moreover we have already proved that they permute the elements of the words, so that their restrictions to the symmetric group $\mathcal{S}_n$ are two bijections.

We can also show the identity of equation~\eqref{eq:des_lec} by an inductive argument on the length $\ell$ of $w\in W$. If $\ell=0$, then $w=\varepsilon$ is the empty word and $\lec(\psi(\varepsilon))=\lec(\varepsilon)=0=\des(\varepsilon$. If $w$ has length $\ell>0$ we have that
\begin{align}
    \lec(\psi(w))&=\lec(\psi(w\setminus\sw(w)))+\lec(\delta(\sw(w))\notag\\
    &=\des(w\setminus\sw(w))+\des(\sw(w))=\des(w),
\end{align}
where the first equality holds since $\delta(\sw(w))$ is the rightmost hook of $\psi(w)$, the second is true by inductive hypothesis and by definition of $\delta$. The last equality is true thanks to Proposition~\ref{prop:des_spezza}.
\end{proof}

\end{document}
